\SourcePage{642}{645}
\section{The scattering amplitude in the momentum representation}

The scattering amplitude involves only the directions of the scattered particle's initial and final momenta. It is therefore natural that the same concept can

\SourcePage{643}{646}
also be reached by formulating the scattering problem in the momentum representation, where the spatial distribution of the process as a whole is not considered at all. Let us show how this is done.


First transform the original Schrödinger equation
\begin{equation}
 -\frac{\hbar^2}{2m}\Delta\psi(\mathbf r)+[U(\mathbf r)-E]\psi(\mathbf r)=0
 \tag{130.1}
\end{equation}
to the momentum representation, passing from coordinate wave functions to
their Fourier components
\begin{equation}
 a(\mathbf q)=\int\psi(\mathbf r)e^{-i\mathbf{qr}}dV.
 \tag{130.2}
\end{equation}
Conversely,
\begin{equation}
 \psi(\mathbf r)=\int a(\mathbf q)e^{i\mathbf{qr}}
 \frac{d^3q}{(2\pi)^3}.
 \tag{130.3}
\end{equation}

Multiply (130.1) by $e^{-i\mathbf{qr}}$ and integrate over $dV$. Two
integrations by parts in the first term give
\[
 \int e^{-i\mathbf{qr}}\Delta\psi(\mathbf r)dV
 =\int\psi(\mathbf r)\Delta e^{-i\mathbf{qr}}dV=-q^2a(\mathbf q).
\]
In the second term, inserting (130.3), we find
\[
 \begin{aligned}
 \int U(\mathbf r)\psi(\mathbf r)e^{-i\mathbf{qr}}dV
 &=\iint U(\mathbf r)e^{-i\mathbf{qr}}a(\mathbf q')
 e^{i\mathbf q'\mathbf r}dV\frac{d^3q'}{(2\pi)^3}\\
 &=\int U(\mathbf q-\mathbf q')a(\mathbf q')\frac{d^3q'}{(2\pi)^3},
 \end{aligned}
\]
where $U(\mathbf q)$ denotes the Fourier component of $U(\mathbf r)$
\footnote{For notational convenience, $\mathbf q$ is written as an argument
of the Fourier component rather than as a subscript.}:
\[
 U(\mathbf q)=\int U(\mathbf r)e^{-i\mathbf{qr}}dV.
\]
Thus the Schrödinger equation in the momentum representation is
\begin{equation}
 \left(\frac{\hbar^2q^2}{2m}-E\right)a(\mathbf q)
 +\int U(\mathbf q-\mathbf q')a(\mathbf q')\frac{d^3q'}{(2\pi)^3}=0.
 \tag{130.4}
\end{equation}

\SourcePage{644}{647}
Notice that this is an integral equation, not a differential equation.

Write the wave function describing the scattering of particles with momentum
$\hbar\mathbf k$ as
\begin{equation}
 \psi_{\mathbf k}(\mathbf r)=e^{i\mathbf{kr}}+\chi_{\mathbf k}(\mathbf r),
 \tag{130.5}
\end{equation}
where $\chi_{\mathbf k}(\mathbf r)$ has asymptotically, as $r\to\infty$, the
form of an outgoing spherical wave. Its Fourier component is
\begin{equation}
 a_{\mathbf k}(\mathbf q)=(2\pi)^3\delta(\mathbf q-\mathbf k)
 +\chi_{\mathbf k}(\mathbf q),
 \tag{130.6}
\end{equation}
and substitution in (130.4) gives the following equation for
$\chi_{\mathbf k}(\mathbf q)$\footnote{By the properties of the
delta-function, the product $(q^2-k^2)\delta(\mathbf q-\mathbf k)$ gives zero
when multiplied by an arbitrary function $f(\mathbf q)$ nonsingular at
$\mathbf q=\mathbf k$ and integrated over $d^3q$. In this sense,
$(q^2-k^2)\delta(\mathbf q-\mathbf k)=0$.}:
\begin{equation}
 -\frac{\hbar^2}{2m}(k^2-q^2)\chi_{\mathbf k}(\mathbf q)
 =U(\mathbf q-\mathbf k)+\int U(\mathbf q-\mathbf q')
 \chi_{\mathbf k}(\mathbf q')\frac{d^3q'}{(2\pi)^3}.
 \tag{130.7}
\end{equation}

It is useful to transform this equation by introducing a new unknown function
in place of $\chi_{\mathbf k}(\mathbf q)$, defined by
\begin{equation}
 \chi_{\mathbf k}(\mathbf q)=-\frac{2m}{\hbar^2}
 \frac{F(\mathbf k,\mathbf q)}{q^2-k^2-i0}.
 \tag{130.8}
\end{equation}
This removes the singularity at $q^2=k^2$ from the coefficients in (130.7),
which becomes
\begin{equation}
 F(\mathbf k,\mathbf q)=U(\mathbf q-\mathbf k)-\frac{2m}{\hbar^2}
 \int\frac{U(\mathbf q-\mathbf q')F(\mathbf k,\mathbf q')}
 {q'^2-k^2-i0}\frac{d^3q'}{(2\pi)^3}.
 \tag{130.9}
\end{equation}
The term $i0$, meaning the limit $i\varepsilon$ as $\varepsilon\to+0$, has
been included in (130.8) to assign a definite meaning to the integral in
(130.9): it specifies how the pole $q'^2=k^2$ is bypassed (compare \S~43). We
shall show that this prescription gives the required asymptotic form of
\begin{equation}
 \chi_{\mathbf k}(\mathbf r)=-\frac{2m}{\hbar^2}\int
 \frac{F(\mathbf k,\mathbf q)e^{i\mathbf{qr}}}{q^2-k^2-i0}
 \frac{d^3q}{(2\pi)^3}.
 \tag{130.10}
\end{equation}

Write $d^3q=q^2dq\,do_q$ and first integrate over $do_q$, the directions of
$\mathbf q$ relative to $\mathbf r$. An integration of this kind was already
carried out in transforming the first term of (125.2); at large $r$ it gives

\SourcePage{645}{648}
\[
 \chi_{\mathbf k}(\mathbf r)=-\frac{2m}{\hbar^2}\frac{2\pi i}{r}
 \int_0^\infty
 \frac{F(\mathbf k,q\mathbf n')e^{iqr}-F(\mathbf k,-q\mathbf n')e^{-iqr}}
 {q^2-k^2-i0}\,q\,dq\frac1{(2\pi)^3},
\]
where $\mathbf n'=\mathbf r/r$, or
\[
 \chi_{\mathbf k}(\mathbf r)=-\frac{im}{2\pi^2\hbar^2r}
 \int_{-\infty}^{\infty}
 \frac{F(\mathbf k,q\mathbf n')e^{iqr}q\,dq}{q^2-k^2-i0}.
\]

The integrand has poles at $q=k+i0$ and $q=-k-i0$, bypassed below and above,
respectively, in integration through the complex $q$ plane (Fig.~48a). Shift
the contour slightly into the upper half-plane, replacing it by a straight
line parallel to the real axis and a closed loop enclosing the pole $q=k$
(Fig.~48b). As $r\to\infty$, the integral along the straight line vanishes
because the integrand contains $\exp(-r\mathop{\rm Im}q)$. The closed-loop
integral is the residue at $q=k$ multiplied by $2\pi i$. We finally obtain
\begin{equation}
 \chi_{\mathbf k}(\mathbf r)=-\frac{m}{2\pi\hbar^2}
 \frac{e^{ikr}}rF(k\mathbf n,k\mathbf n'),
 \tag{130.11}
\end{equation}
where $\mathbf n$ is the unit vector along $\mathbf k$. Thus the required
asymptotic form is obtained, and the scattering amplitude is
\begin{equation}
 f(\mathbf n,\mathbf n')=-\frac{m}{2\pi\hbar^2}
 F(k\mathbf n,k\mathbf n').
 \tag{130.12}
\end{equation}
The scattering amplitude is therefore determined by the value at $q=k$ of the
function $F(\mathbf k,\mathbf q)$ satisfying (130.9).

\begin{figure}[htbp]
\centering
\begin{tikzpicture}[x=1cm,y=1cm,>=stealth,line width=.5pt,font=\small]
 \begin{scope}[shift={(-2.5,0)}]
  \draw[->] (-1.7,0)--(-.45,0); \draw (-.45,0) arc[start angle=180,end angle=0,radius=.25];
  \draw[->] (.05,0)--(1.05,0); \draw (1.05,0) arc[start angle=180,end angle=360,radius=.25];
  \draw[->] (1.55,0)--(1.85,0);
  \fill (-.2,0) circle (1.3pt); \fill (1.3,0) circle (1.3pt);
  \node[circle,draw,inner sep=2pt] at (-1.55,.8) {$q$}; \node[below] at (0,-.55) {(a)};
 \end{scope}
 \begin{scope}[shift={(2.3,0)}]
  \draw (-1.5,0)--(1.5,0); \fill (0,0) circle (2pt);
  \draw[->] (-1.5,.55)--(-.65,.55); \draw (-.65,.55)--(.2,.55);
  \draw[->] (.2,.55)--(.95,.55); \draw (.95,.55)--(1.5,.55);
  \draw[->] (0,.55) arc[start angle=90,end angle=-270,radius=.3];
  \node[below] at (0,-.55) {(b)};
 \end{scope}
\end{tikzpicture}
\caption{}
\label{fig:48}
\end{figure}

When perturbation theory applies, (130.9) is readily solved by successive
iteration. In the first approximation, omitting the integral term altogether,

\SourcePage{646}{649}
one has $F(\mathbf k,\mathbf q)=U(\mathbf q-\mathbf k)$. In the next
approximation, insert this first-order expression for $F$ in the integral term.
After a slight change of variables, (130.12) then gives
\begin{equation}
 \begin{split}
 f(\mathbf n,\mathbf n')=-\frac{m}{2\pi\hbar^2}\Bigg\{
 &U(\mathbf k'-\mathbf k)\\
 &+\frac{2m}{\hbar^2}\int
 \frac{U(\mathbf k'-\mathbf k'')U(\mathbf k''-\mathbf k)}
 {k^2-k''^2+i0}\frac{d^3k''}{(2\pi)^3}\Bigg\},
 \end{split}
 \tag{130.13}
\end{equation}
where $\mathbf k=k\mathbf n$ and $\mathbf k'=k\mathbf n'$. The first term is
the first-Born-approximation formula (126.4); the second is the contribution of
the second approximation to the scattering amplitude\footnote{This result can
of course also be obtained without passing to the momentum representation.
Comparison of (43.1) and (43.6) makes it evident that the second-approximation
formula differs from the first by replacing $U(\mathbf k'-\mathbf k)$ with the
expression in braces in (130.13).}.

Formula (130.13) demonstrates the fact mentioned in \S~126: already in the
second approximation the scattering amplitude loses the symmetry (126.8). At
first sight the integral term in (130.13) might also seem symmetric under
interchange of the initial and final states. It is not, because complex
conjugation changes the integration contour, namely the direction in which the
pole is bypassed.
